\scp{Reviewing CMP}{15-02}
Simply put,
we find no support for CMP in the historical phonology or the morphosyntax.
We do not know what defines a CMP language,
nor what distinguishes them from languages that are not CMP.
One of our proposals in this book is that CMP does not exist,
and instead it should be replaced with the nine subgroups we
identify -- one of which links historically with \tsc{Celebic} languages (\crf{ch:Tal}).
\ili{As} yet, we have found no evidence to unify the eight Wallacean
subgroups in this region together at a node lower than MP
(see \frf{01-06}, \prf{fig:01-06}).

For example, the data in the region do not support the merger
of “PAn *-mb- and *-mp- as PCMP *mb” \citep[26]{Blust1981}.
The data in the region do not support the “post-nasal voicing
of stops,” seen as “a highly distinctive development not known
to be shared with any language outside eastern Indonesia” to
be a feature supporting PCMP (see \cite[94--98]{Blust2008}, \cite[72]{Blust2009}).
Responding to criticisms, \citet[72]{Blust2009} conceded this
feature “was a product of diffusion” (see discussion in \srf{13-01-02}).
The loss of antepenultimate vowels,
particularly in PMP trisyllables beginning with *ha-
and *qa- \citep[263f]{Blust1993},
also cannot be taken as diagnostic of CMP.
\citet[123]{DonohueGrimes2008} pointed out sporadic retentions
in several etyma throughout the region.
For instance, languages as far \ili{east} as
\ili{Arguni} (STB, Papua) retain the penultimate
vowel of *qalu-hipan > **qalipan \it{arifen} `centipede',
as does \ili{Hoti} of eastern Seram with \it{urial} `centipede'.

In \srf{11-03-01} we noted that \ili{Watubela} preserves PMP *q in
all positions, including in trisyllabic words like PMP *qatəluR > \it{katlu} ‘egg’.
And in \srf{10-03-01} we noted the regular sound correspondences
and conditioning environments for Boano explained the different
reflexes of PMP *limas > lima-te ‘bailer (n)’ versus *qalima > \it{nima} ‘hand’.
In Boano, initial *l = \it{l},
but medial (intervocalic) *l > \it{n}.
These kinds of things are trace evidence for PMP *qa-
and antepenultimate vowels in the region.
This same phenomenon is also found in languages west of the Blust line.

\citet[264ff]{Blust1993} sees glide truncation (the loss of the coda
in historical vowel-glide sequences) as “one of the most distinctive
characteristics of the phonological history of CMP languages”.
In \srf{13-01-03} we showed that several vowel-glide sequences
are preserved in full in some languages
and in several of our subgroups,
and in others the variety of reflexes require that the full historical
vowel-glide sequence be reconstructed for those as well.
Blust’s “glide truncation” does not characterise any of our Wallacean subgroups.
In response to criticisms from \citet{DonohueGrimes2008},
\citet[72]{Blust2009} conceded that “glide truncation,
was a product of diffusion,
since some languages that show no evidence of it are interspersed
with those that do.”

\newpage
\largerpage
That leaves only the lexicon as the main basis for a major subgroup
in a region known for its intense contact with Papuan languages
spanning almost four millennia (see \srf{02-03-01}, \srf{15-Features}, and \srf{15-04}),
and for on-going world-impacting trade spanning
at l\ili{east} two millennia (see \srf{02-02-01}).
Methodologically the lexicon is not only the weakest basis for
subgrouping, but also the most problematic.
Furthermore, if there is no CMP,
then there is also no proto-language to which to assign lexical reconstructions.
It is our view that many PCMP etyma that have been reconstructed
must be seen as spurious.
They do not add insight into the patterned development of the
lexicon or phonology or morphology.
They do not account for the many languages in the region that
do not reflect that PCMP reconstruction
and better reflect the PMP form,
or something more localised.
Many developments are better explained directly through normal
developments in the historical phonology from the PMP forms,
such as: regular loss of some final consonants in some languages
(particularly on verbs), such as PMP *balabuR ‘blurred,
of vision’ > PCMP **balabu ‘dim,
of vision’; frozen stative or causative prefixes for some languages
that are still productive in others,
and not present in still others; widespread weakening of antepenultimate
vowels that are still reflected in some languages in some items;
widespread loss of PMP *h
and *q; widespread but sporadic coalescence of secondary vowel
sequences through the loss of historical consonants, and so forth.
All PCMP reconstructions are in need of review.

Throughout this work some reflexes have been identified from
languages in the target region that are actually consistent with
forms that have been reconstructed for Proto-Philippines
and PWMP, and should probably be reassigned to PMP.

There are also a few etyma that are found in two or more of our
Wallacean subgroups, but not throughout all of them.
Many of these lower-level etyma with limited distribution have
been flagged with a double asterisk ** throughout this work
and are identified at various levels.
These etyma raise methodological questions.
The forms in daughter languages in different subgroups often
reflect the same regular sound correspondences we find in those languages from PMP.
But their distribution is sporadic
and limited, and often only reflected in certain nodes of some
subgroups and not others (much like {\#}muku ‘banana’, see \srf{02-03-01-07}).
Some are distributed along (parts of) the southern string of
islands, and others in (parts of) the central Maluku region.
There is no common parent language based on the historical phonology
to which to assign these etyma-of-convenience,
so they are best seen as words that reflect diffusion through
contact -- which is known to be prevelant in the region.
They are useful for understanding the lexicon,
but they are not useful for subgrouping.
Words of this sort are discussed further in \srf{15-04}.